\section{Metodologia}
\label{sec:metodologia}

\pendente{Descrever o percurso metodológico completo, incluindo classificação
da pesquisa, procedimentos de levantamento de requisitos, processo de
desenvolvimento e forma de avaliação.}

\subsection{Classificação da pesquisa}

\pendente{Natureza, objetivos e procedimentos técnicos, com a fonte
metodológica que sustenta a classificação.}

\subsection{Levantamento de requisitos}

\pendente{Entrevistas com servidores do setor de licitações e análise de
documentos reais de processos já encerrados. Registrar em
\code{TCC/requisitos/} cada sessão: o que foi dito, o que foi inferido e o que
ficou em aberto.}

\subsection{Tratamento dos dados de domínio}

Documentos reais de processos licitatórios não são versionados nem reproduzidos
no texto. Cada documento recebido é mantido fora do repositório e dele se deriva
uma versão sintética, com a mesma estrutura e nenhum dado verdadeiro, que é a
única utilizada em testes, exemplos e ilustrações deste trabalho.

\pendente{Detalhar o procedimento de derivação e discutir suas limitações para
a avaliação da solução.}

\subsection{Desenho da solução}

O desenho que resultou do levantamento --- requisitos, regras de negócio e
arquitetura --- é apresentado na seção~\ref{sec:requisitos-arquitetura}, e não
aqui. A separação é deliberada: ele é resultado do percurso metodológico e não
etapa dele, e tratá-lo como procedimento esconderia que boa parte das escolhas
foi determinada por evidência de campo, não por preferência técnica.

Duas escolhas, porém, pertencem à metodologia, porque dizem respeito à natureza
da evidência que as sustenta.

\subsubsection{Recuperação de informação sobre a legislação e os modelos}

\pendente{Recuperação aumentada por geração aplicada à legislação e aos modelos
de documento já aprovados. Descrever a estratégia de fragmentação, o modelo de
representação escolhido e o critério de recuperação. Há material reaproveitável
dos fichamentos da Iniciação Científica, mas cada citação precisa ser
reconferida contra o PDF.}

\pendente{Posicionar contra o contra-argumento levantado no levantamento de
literatura: há trabalho publicado sustentando que recuperação plana sobre
documento não estruturado limita raciocínio de múltiplos saltos quando os dados
são estruturados. Responder explicitamente, não por omissão.}

\subsubsection{Natureza da evidência que sustenta a escolha da persistência}

A separação entre banco relacional e índice semântico, descrita na
subseção~\ref{subsec:arquitetura}, apoia-se em medição realizada pelo autor em
iniciação científica anterior a este trabalho, e não em literatura revisada por
pares. O texto registra isso onde a medição é invocada, e a distinção é mantida
deliberadamente: medição própria é evidência, com força menor que a de resultado
publicado e submetido a revisão.

\pendente{Descrever o protocolo daquela medição --- conjunto de dados,
seletividade do filtro, métrica de recuperação e número de repetições --- em
detalhe suficiente para que o leitor julgue se ela transfere para o caso deste
trabalho. Critério acrescentado pela orientação: verificar contra a natureza dos
dados a tratar, não só o desempenho medido.}

\subsection{Desenvolvimento}

\pendente{Processo de desenvolvimento, com testes automatizados como critério
de conclusão de cada incremento.}

\subsection{Avaliação}

\pendente{Definir como a solução será avaliada e contra o quê. Sem uma linha de
base medida antes da intervenção, não há como afirmar redução de esforço ---
planejar a coleta dessa linha de base ainda em TCC I.}

\pendente{A etapa de verificação pelo órgão de controle ficou fora do escopo
obrigatório por decisão da orientação em 19/08/2026: o tempo de resposta do
órgão não é previsível e não cabe em planejamento. Reaproveitável no fim, se
houver tempo.}
